% Pure results fragment for the existing ctexart report. No preamble/document wrapper.
% Sources: structured_information_results_review.md/.json;
% evidence/structured_information/v3_frozen/{calibration,evaluation}/summary.json;
% v3_frozen/stress_v_correlation/summary.json and failure_diagnostic/summary.json.
% Original protocol primary endpoint: paired continuous signed margin.
% This fragment has not been compiled or opened, as requested.

\subsection{冻结正式批次：39条校准历史与80条评价历史}
\label{sec:information-frozen-results}
正式批次完成39条独立校准历史和80条独立评价历史；评价中的79条历史完成全部11种方法，另1条保留恒充的联合拟合失败。信息设计的80条评价均完成。评价真结构为star 10条、weak 12条、moderate 29条、strong 29条；公开计量质量为uniform 22条、temporal 27条、spatial 31条。每条历史的11种方法共享历史读数、公开质量、真实背景和标准化未来噪声，校准与评价各历史只抽取一种公开质量。四种star/cherry候选都进入最终拟合，候选漏真为0/80。全部119条记录的策略标识和19份科学源码哈希与冻结记录一致。

主要结果是：信息设计相对恒充提高了实际联合拟合裕量，但同时增加了误拒真候选的次数；相对简单公开质量规则，其全总体平均裕量优势尚未出现。为保持原协议的评价顺序，主要终点是同一历史下的连续裕量差；本结果部分的 $n_H$ 表示相应的有效历史数。令
\begin{equation}
 M=\min_{h\ne h_{\rm true}}\mathrm{SSE}_h-\mathrm{SSE}_{h_{\rm true}},
 \label{eq:information-actual-margin-v3}
\end{equation}
其中$\mathrm{SSE}_h$是候选$h$用实际历史和新增计量联合拟合取得的完整白化平方误差。$M>0$表示真候选的损失较小；在所有候选拟合完成且真候选被保留时，$M$超过门槛才产生唯一正确结果。真结构只用于这个离线评价量。部署选择仍仅使用历史及公开质量。

\subsubsection{互斥结果与主要配对比较}
表\ref{tab:information-outcomes-v3}将结果分为四类：唯一正确、安全未决、误拒真候选、计算或物理失败。安全未决要求真候选仍被保留且至少还保留一个错误候选。每行四类之和为80。旧字段\texttt{unresolved}只检查保留候选数是否大于1，会与误拒和失败重叠，因此本表采用独立审查重算的互斥计数。连续裕量只用完成联合拟合的历史，失败历史保留在二元结果分母中，不补填裕量。

{\small
\begin{longtable}{L{.31\textwidth}rrrrr}
\caption{正式80条评价历史的互斥结果；最后一列为$M$均值及有效历史数。}\label{tab:information-outcomes-v3}\\
\toprule
方法 & 唯一正确 & 安全未决 & 误拒 & 失败 & 平均$M$（$n_H$）\\
\midrule\endfirsthead
\toprule
方法 & 唯一正确 & 安全未决 & 误拒 & 失败 & 平均$M$（$n_H$）\\
\midrule\endhead
无EV（none） &28&48&4&0&4.670（80）\\
固定恒充（constant） &49&27&3&1&225.704（79）\\
简单全局质量规则 &53&22&5&0&917.562（80）\\
简单本地质量规则 &52&23&5&0&879.411（80）\\
固定节点信息优化 &53&23&4&0&643.972（80）\\
随机节点信息优化 &51&24&5&0&717.241（80）\\
最佳操作初值（seed） &55&20&5&0&865.293（80）\\
信息设计（designed） &57&16&7&0&887.805（80）\\
单中心（nominal） &57&16&7&0&891.309（80）\\
等权任务（uniform） &56&17&7&0&873.820（80）\\
旧profile策略 &48&26&6&0&595.303（80）\\
\bottomrule
\end{longtable}}

信息设计相对恒充的平均配对裕量差为$627.891$，配对标准误为$127.681$，有效分母79；差值中位数为$113.659$，53条提高、16条下降、10条在$10^{-6}$容差内相同。这一比较对应预声明的主要方法比较。唯一正确则由49/80增加为57/80；全80条配对中9胜、1负，其余70条两方法结果相同。其中1胜来自恒充拟合失败，而非已完成拟合之间的额外排错。双侧精确二元配对$p=0.0215$作为补充描述；它不替代原协议的连续裕量主终点。

{\small
\begin{longtable}{L{.27\textwidth}rrrr}
\caption{信息设计减去各对照的同历史配对结果。SE为配对均值差的标准误；唯一正确胜/负始终采用全80条历史。}\label{tab:information-paired-v3}\\
\toprule
对照 & 有效$n_H$ & 平均$M$差$\pm$SE & 差中位数 & 唯一正确胜/负\\
\midrule\endfirsthead
\toprule
对照 & 有效$n_H$ & 平均$M$差$\pm$SE & 差中位数 & 唯一正确胜/负\\
\midrule\endhead
无EV &80&$883.135\pm148.740$&187.259&31/2\\
固定恒充 &79&$627.891\pm127.681$&113.659&9/1\\
简单全局质量规则 &80&$-29.757\pm147.447$&0.000&6/2\\
简单本地质量规则 &80&$8.394\pm146.016$&0.000&7/2\\
固定节点信息优化 &80&$243.833\pm102.640$&0.000&6/2\\
随机节点信息优化 &80&$170.564\pm69.889$&0.000&10/4\\
最佳操作初值 &80&$22.512\pm12.353$&0.000&4/2\\
单中心 &80&$-3.504\pm3.509$&0.000&0/0\\
等权任务 &80&$13.985\pm57.470$&0.000&4/3\\
旧profile策略 &80&$292.502\pm130.536$&0.955&11/2\\
\bottomrule
\end{longtable}}

简单全局规则的平均裕量较高，而信息设计唯一正确数较多，说明提高已清楚案例的裕量与使新案例排除全部错误候选是两种收益。信息设计相对全局规则的配对均值为$-29.757$，SE为$147.447$，唯一正确配对6胜/2负；相对本地规则均值为$8.394$，SE为$146.016$，7胜/2负。多中心与单中心的唯一正确结果逐历史相同，77/80条裕量差在$10^{-6}$内相同。后续搜索相对最佳初值平均增益$22.512$，但52/80条裕量相同；任务加权相对等权为4胜/3负。这些消融保留了各项复杂化的实际增益大小。

\subsubsection{star、弱边与历史启用范围}
信息设计在moderate中27/29条唯一正确，在strong中28/29条唯一正确；在star中为0/10且误拒4/10，在weak中为2/12且误拒2/12。star与weak的所有必要拟合都完成，因此这些结果属于辨识或门控结果。表\ref{tab:information-case-v3}保留四类结构，strong的大裕量增益与新增唯一正确数应分别读取。

{\small
\begin{longtable}{lrrrr}
\caption{预声明结构分层。四元组依次为唯一正确/安全未决/误拒/失败；裕量差为信息设计减恒充。}\label{tab:information-case-v3}\\
\toprule
结构 & $n_H$ & 信息设计 & 恒充 & 平均$M$差（有效$n_H$）\\
\midrule\endhead
star &10&0/6/4/0&0/9/1/0&$-2.315$（10）\\
weak &12&2/8/2/0&0/11/1/0&0.196（12）\\
moderate &29&27/1/1/0&22/6/1/0&139.724（29）\\
strong &29&28/1/0/0&27/1/0/1&1627.578（28）\\
\bottomrule
\end{longtable}}

按历史选择状态和历史中心状态同时启用的规则，实际信息搜索启用70条历史：star 10、weak 12、moderate 29、strong 19。该子组信息设计47/70条唯一正确，恒充39/70，全局/本地质量规则43/70和42/70。信息设计相对全局规则的平均配对裕量差为$282.378$（SE $116.968$），相对本地为$315.232$（SE $116.458$）。分组规则只读历史，没有按未来结果挑选案例。

另10条均为strong：8条的历史间隔超过启用门槛，2条的star候选超出历史计算支撑范围，信息策略按冻结规则回退恒充。这10条中上述四种策略都10/10唯一正确，但全局质量规则的平均裕量为3005.781，回退恒充为791.076。该子组解释了全局规则为何在全80平均裕量上更高，却没有为其增加该组唯一正确数。总体比较仍使用全80条；启用子组的结果另列，整体校准并未提供该子组的条件风险结论。

\subsubsection{校准门槛、误拒与同历史代理检验}
39条校准历史均完成。对每条历史，先取11种方法中真候选相对最优候选的最大完整SSE差，再取39个统计量的第36顺序统计量，得到冻结门槛
\[
 \kappa_\alpha=2.9617106149485366.
\]
在声明混合、独立且交换的历史条件下，有限秩结论给出边际误拒上界0.10，其概率同时平均随机校准数据和一条新历史。当前已经实现的这个固定门槛，在80条评价中任一方法误拒11/80，即13.75\%；Clopper--Pearson 95\%区间为[7.07\%,23.27\%]。区间包含10\%，但当前固定门槛的条件风险仍需更多独立评价来估计。70条启用历史中，任一方法误拒11/70，信息设计误拒7/70，应与整体边际结论分开报告。

代理检查使用同一历史内两个计划的分数差与实际$M$差。恒充和固定节点的代理分数来自保存的robust节点记录，seed与designed使用相同中心、任务权重及尺度。信息设计相对恒充的69条可比历史中，Pearson相关为0.492、Spearman相关为0.455；代理增益全部为正，实际裕量53条提高、16条下降。固定节点相对恒充的相关为0.185/0.184，实际51条提高、18条下降。这里缺少的11条是10条未启用历史及恒充失败的950052。

其余计划比较有70条同历史配对：designed相对fixed的Pearson/Spearman为0.646/0.433，相对seed为0.391/0.288，相对random为0.576/0.439。后续迭代只在25/70条启用历史上提高最终代理分数，其中实际裕量17条提高、8条下降；210个节点搜索中50个节点优于最佳操作初值。代理能提供部分方向信息，逐历史实际增益仍需联合AC拟合核验；star和weak的代理与实际增益联系尤其弱。相关性与多项对照均作描述，未进行多重比较校正。

\subsubsection{原拟合失败的单独恢复诊断}
950052是strong/spatial历史，原恒充对clade\_12、star及clade\_02的联合拟合在每初值200次函数评价处停止。另存诊断直接重建保存的历史和新增读数，使用实际保存的joint计量模型及完全一致的候选支撑顺序。每候选分别用历史warm初值和计量cold初值，各4个起点、每起点上限800；全部24个起点以目标变化容差结束，实际178--236次评价，恢复3/3，无需1500次后备。

三个错误候选的最小已取得局部SSE依次为974.783519、974.776766和974.895474；原真候选clade\_01独立重放为26.248613，诊断裕量948.528153。恢复状态触及若干R/X下界及末端1的X上界0.1 pu。58个保存初值、返回状态和原状态经独立分支扫流重放，SSE、Gaussian对数似然、AC方程及功率平衡全部通过；最大平衡残差为$2.885\times10^{-15}$。诊断共6次拟合、5034次实际AC调用、5029次函数评价及138.839秒，没有失败的AC调用。19份科学源码与138份原证据前后字节一致。

该证据将原950052的停止原因定位为冻结拟合预算问题；原历史仍为未完整完成，恒充仍保留\texttt{joint\_fit\_failure}并计入失败分母，校准门槛和正式表均保持原值。恢复状态提供已取得局部目标，不将其转写为原冻结策略的成功。

\subsubsection{相关重复计量压力与计算成本}
另取正式评价的前16条历史作配对压力检查，沿用已选动作、背景、历史和名义门槛，只使同点电压重复读数之间的相关系数为0.25。若同点重复次数为$m_{t,j}$，实际均值方差相对独立假设的比例为$1+0.25(m_{t,j}-1)$；$m_{t,j}=1,5,9$分别为1、2、3。计量拟合仍使用原独立平均方差，从而直接检查过高精度假设的后果。

16条压力历史全部完成，新增独立历史数为0；这是同历史配对检查。任一方法误拒由原同16条的2/16变为3/16。信息设计的唯一正确/安全未决/误拒为12/2/2，与原同16条相同；全局和本地质量规则均为13/2/1，恒充为12/3/1。无EV误拒由1/16增至3/16，固定和随机节点信息方法各由1/16增至2/16。信息设计的平均压力减名义裕量为0.140，简单全局/本地规则为$-4.510$和$-6.190$。压力批次用了348次实际联合候选拟合、9934次联合AC调用、291.160秒；全部11方法与缓存覆盖176个逻辑动作。名义混合校准风险适用原独立计量模型，压力批次单列。

{\small
\begin{longtable}{L{.48\textwidth}rrr}
\caption{正式评价的保存计时组件，均值分母80；AC列为平均算法尝试次数。}\label{tab:information-cost-v3}\\
\toprule
计算范围 & 均值秒 & 最大秒 & 平均AC调用\\
\midrule\endhead
robust初始化及搜索 &0.722&0.945&---\\
加历史中心构造的robust规划 &0.778&1.009&564.6\\
H拟合、中心、robust及本方法联合拟合组件和 &1.888&2.335&794.4\\
全11方法选择 &15.238&31.292&3901.9\\
全11方法完整benchmark &19.706&36.902&4687.4\\
其中旧profile搜索 &12.569&28.035&2088.2\\
\bottomrule
\end{longtable}}

四种信息搜索含初始化的最大合计3.138秒，全选择最大31.292秒，均通过预声明15/75秒预算门槛。信息设计未缓存联合拟合平均0.661秒、每次4个候选；robust规划连同中心构造平均增加1.75次历史拟合。简单全局/本地质量规则每历史三个操作子问题的求解器时间合计为0.0040/0.0038秒，选择阶段不调用AC信息计算。该solver时间与含Python构造和初始化的robust时间口径不同，不能直接相除作为端到端加速比。完整80条benchmark累计1576.512秒及374988次AC调用。

成本表的组件和由保存记录相加，尚未另跑端到端信息策略；它未含共享候选生成、任务权重和少量调度开销。全11方法benchmark包含所有消融及旧profile，并复用相同动作的联合拟合缓存，因此逐方法记录耗时也不能作为独立部署耗时。旧profile的预测中心与迭代预算不同，比较对象是完整旧策略。

本阶段通过了冻结批次、保存证据重放和计算预算验收；主要连续裕量结果支持该合成混合下信息设计相对恒充的增益。对简单质量规则的额外收益、固定门槛的条件风险，以及多中心、加权任务和后续迭代各自的必要收益，需按本节的配对结果继续判断。所有最终损失来自四个固定候选、同一稳定AC分支上的多初值局部拟合，尚无全局损失下界；操作MILP的界只覆盖对应仿射子问题。上述结论的范围是冻结三末端模型、指定参数与计量混合及原拟合预算，现场模型和其他候选结构另需验证。
